\section{Ekonometrické modely pro časové řady a panelová data}

\subsection*{Základní rámec}

Ekonometrické modely pro \textbf{časové řady} a \textbf{panelová data} rozšiřují klasickou regresi o situace,
kdy pozorování nejsou jednoduše nezávislá. U časových řad záleží na pořadí pozorování v čase,
protože minulost často ovlivňuje současnost. U panelových dat sledujeme stejné jednotky opakovaně
v čase, a proto můžeme kontrolovat nepozorované vlastnosti jednotek, které se v čase nemění.

\begin{itemize}
    \item \textbf{Časová řada:} jedna proměnná nebo více proměnných pozorovaných v čase.
    \[
        y=\{y_t,\ t\in T\}.
    \]
    Typicky řešíme trend, sezónnost, autokorelaci, nestacionaritu a predikci.

    \item \textbf{Panelová data:} stejné jednotky $i=1,\dots,N$ pozorované v čase $t=1,\dots,T$.
    \[
        y_{it}, \qquad i=1,\dots,N,\quad t=1,\dots,T.
    \]
    Typicky řešíme nepozorovanou heterogenitu jednotek, fixed effects, random effects a dynamické panely.
\end{itemize}

\subsection*{1. Časové řady}

\subsubsection*{1.1 Co je časová řada}

Časová řada je posloupnost hodnot pozorovaných v pravidelných časových intervalech, například denně,
měsíčně, kvartálně nebo ročně. Důležité je, že pozorování nelze libovolně přeházet. Hodnota HDP v roce
2024 logicky navazuje na hodnotu z roku 2023, cena akcie dnes souvisí s cenou včera a nezaměstnanost
v jednom měsíci souvisí s nezaměstnaností v předchozích měsících.

\[
    y_t,\qquad t=1,\dots,T.
\]

Hlavní cíle analýzy časových řad jsou:
\begin{itemize}
    \item \textbf{pochopení mechanismu} generujícího data,
    \item \textbf{predikce budoucích hodnot},
    \item \textbf{testování hypotéz} o dynamice procesu,
    \item \textbf{modelování vztahů mezi proměnnými v čase}.
\end{itemize}

Oproti průřezovým datům je problém v tom, že pozorování v čase často nejsou nezávislá:
\[
    y_t \text{ může záviset na } y_{t-1}, y_{t-2}, \dots
\]

\subsubsection*{1.2 Komponenty časové řady}

Časovou řadu si často představujeme jako kombinaci několika složek:
\[
    y_t = Tr_t + S_t + C_t + u_t,
\]
kde:
\begin{itemize}
    \item $Tr_t$ je trend,
    \item $S_t$ je sezónní složka,
    \item $C_t$ je cyklická složka,
    \item $u_t$ je náhodná složka.
\end{itemize}

V jednodušší podobě:
\[
    y_t = Tr_t + u_t.
\]

\paragraph{Trend}

Trend znamená dlouhodobý růst nebo pokles časové řady. Může být deterministický nebo stochastický.

\textbf{Lineární deterministický trend:}
\[
    y_t = \alpha + \beta t + u_t.
\]

\textbf{Kvadratický trend:}
\[
    y_t = \alpha + \beta t + \gamma t^2 + u_t.
\]

Pokud trend odhadneme, detrendovaná řada je:
\[
    \widetilde y_t = y_t - \widehat{Tr}_t.
\]

Intuice: pokud chceme zkoumat vztah dvou ekonomických proměnných, neměli bychom ho zaměňovat
za fakt, že obě proměnné pouze rostou v čase. Například počet televizních seriálů i nějaký společenský
ukazatel mohou růst současně, ale to samo o sobě neznamená ekonomickou souvislost.

\paragraph{Zdánlivá regrese}

U nestacionárních časových řad může vzniknout \textbf{spurious regression}, tedy zdánlivá regrese.
Regrese může mít vysoké $R^2$ a statisticky významné koeficienty, přestože mezi proměnnými není
skutečný vztah. Často je důvodem společný trend.

Typický příklad:
\[
    y_t = \alpha + \beta x_t + u_t,
\]
kde $y_t$ i $x_t$ mají trend. Model pak může zachycovat hlavně to, že obě řady rostou s časem.

Řešení:
\begin{itemize}
    \item zahrnout trend $t$ do modelu,
    \item použít detrendovaná data,
    \item diferencovat časovou řadu,
    \item testovat jednotkový kořen a kointegraci.
\end{itemize}

\paragraph{Sezónnost}

Sezónnost znamená pravidelně se opakující vzorec v rámci roku, týdne nebo jiného cyklu.
Například maloobchodní tržby bývají vyšší v prosinci, spotřeba elektřiny se liší mezi zimou a létem.

Sezónnost lze modelovat pomocí dummy proměnných:
\[
    y_t =
    \alpha+\beta t+\delta_2D_{2t}+\cdots+\delta_sD_{st}+u_t.
\]

Pokud máme například měsíční data, používáme typicky 11 dummy proměnných, protože jedna sezóna
slouží jako referenční kategorie. Použití všech 12 dummy proměnných spolu s konstantou by vedlo
k perfektní multikolinearitě.

\subsubsection*{1.3 Stacionarita}

Stacionarita je jedna z nejdůležitějších vlastností časových řad. Obecně znamená, že základní
pravděpodobnostní vlastnosti procesu se v čase nemění. Pro ekonometrické modelování je důležitá
hlavně \textbf{slabá stacionarita}.

Proces $\{y_t\}$ je slabě stacionární, pokud:
\[
    E(y_t)=\mu
\]
je konstantní pro všechna $t$,

\[
    Var(y_t)=\gamma(0)
\]
je konstantní a

\[
    Cov(y_t,y_{t-h})=\gamma(h)
\]
závisí pouze na délce zpoždění $h$, nikoliv na konkrétním čase $t$.

Jinými slovy:
\begin{itemize}
    \item řada nemá měnící se průměr,
    \item řada nemá měnící se varianci,
    \item vztah mezi současností a minulostí závisí pouze na vzdálenosti v čase.
\end{itemize}

\paragraph{Bílý šum}

Bílý šum je základní stavební kámen mnoha modelů časových řad:
\[
    e_t \sim WN(0,\sigma^2).
\]

Platí:
\[
    E(e_t)=0,
\]
\[
    Var(e_t)=\sigma^2,
\]
\[
    Cov(e_t,e_s)=0 \quad \text{pro } t\neq s.
\]

Bílý šum je tedy posloupnost náhodných chyb bez autokorelace.

\subsubsection*{1.4 Autokorelace, ACF a PACF}

Autokorelace znamená korelaci časové řady se svou vlastní minulostí.

Autokovariance:
\[
    \gamma(h)=Cov(y_t,y_{t-h}).
\]

Autokorelace:
\[
    \rho(h)=\frac{\gamma(h)}{\gamma(0)}.
\]

\paragraph{ACF}

ACF, tedy autocorrelation function, ukazuje korelaci mezi $y_t$ a $y_{t-h}$ pro různá zpoždění $h$.
Používá se k odhalení persistence, sezónnosti a nestacionarity.

U stacionární řady ACF obvykle rychle klesá k nule. U nestacionární řady ACF často klesá velmi pomalu.

\paragraph{PACF}

PACF, tedy partial autocorrelation function, měří vztah mezi $y_t$ a $y_{t-h}$ po očištění vlivu
mezilehlých zpoždění $y_{t-1},\dots,y_{t-h+1}$.

\paragraph{Použití ACF a PACF při volbě modelu}

\begin{center}
\begin{tabular}{|p{0.35\textwidth}|p{0.55\textwidth}|}
\hline
\textbf{Chování ACF/PACF} & \textbf{Typický model} \\
\hline
PACF se po určitém zpoždění utne, ACF postupně doznívá & AR($p$) \\
\hline
ACF se po určitém zpoždění utne, PACF postupně doznívá & MA($q$) \\
\hline
ACF i PACF postupně doznívají & ARMA($p,q$) \\
\hline
ACF klesá velmi pomalu & podezření na nestacionaritu nebo jednotkový kořen \\
\hline
\end{tabular}
\end{center}

V praxi se ACF/PACF používají společně s informačními kritérii, jako jsou AIC a BIC, a s kontrolou reziduí.

\subsubsection*{1.5 Základní lineární modely časových řad}

\paragraph{MA(q) model}

Moving average model řádu $q$:
\[
    y_t = \alpha + e_t+\theta_1e_{t-1}+\cdots+\theta_qe_{t-q}.
\]

Současná hodnota závisí na současném a minulých náhodných šocích.

\paragraph{AR(p) model}

Autoregresní model řádu $p$:
\[
    y_t = \alpha+\phi_1y_{t-1}+\cdots+\phi_py_{t-p}+e_t.
\]

Současná hodnota závisí na vlastních minulých hodnotách. Například AR(1):
\[
    y_t=\alpha+\phi y_{t-1}+e_t.
\]

AR(1) je stacionární, pokud:
\[
    |\phi|<1.
\]

Pokud $\phi$ leží blízko jedné, proces je velmi perzistentní. Pokud $\phi=1$, dostáváme jednotkový
kořen a proces je nestacionární.

\paragraph{ARMA(p,q) model}

ARMA model kombinuje autoregresní a moving average složku:
\[
    y_t =
    \alpha+
    \sum_{i=1}^{p}\phi_i y_{t-i}
    +
    \sum_{j=1}^{q}\theta_j e_{t-j}
    +e_t.
\]

Používá se pro stacionární časové řady. Pokud řada není stacionární, je nutné ji nejprve transformovat,
například diferencováním.

\paragraph{ARIMA(p,d,q) model}

ARIMA model je rozšířením ARMA modelu pro nestacionární řady, které se stanou stacionárními po
diferencování.

První diference:
\[
    \Delta y_t = y_t-y_{t-1}.
\]

Druhá diference:
\[
    \Delta^2 y_t = \Delta y_t-\Delta y_{t-1}.
\]

ARIMA model:
\[
    \Delta^d y_t =
    \alpha+
    \sum_{i=1}^{p}\phi_i \Delta^d y_{t-i}
    +
    \sum_{j=1}^{q}\theta_j e_{t-j}
    +e_t.
\]

Parametry:
\begin{itemize}
    \item $p$ = počet autoregresních zpoždění,
    \item $d$ = počet diferencování nutných ke stacionaritě,
    \item $q$ = počet zpožděných šoků.
\end{itemize}

Speciální případ:
\[
    y_t=y_{t-1}+e_t
\]
je random walk, tedy ARIMA$(0,1,0)$.

\paragraph{SARIMA}

Pokud má řada sezónnost, používá se SARIMA:
\[
    SARIMA(p,d,q)\times(P,D,Q)_s,
\]
kde:
\begin{itemize}
    \item $s$ je délka sezóny,
    \item $P,D,Q$ jsou sezónní analogie parametrů $p,d,q$.
\end{itemize}

Například pro měsíční data bývá:
\[
    s=12.
\]

\subsubsection*{1.6 Jednotkový kořen a ADF test}

Jednotkový kořen znamená, že časová řada je nestacionární a šoky mají trvalý efekt. Typickým příkladem
je random walk:
\[
    y_t=y_{t-1}+e_t.
\]

Základní AR(1) proces:
\[
    y_t=\alpha+\phi y_{t-1}+e_t.
\]

Odečteme $y_{t-1}$:
\[
    y_t-y_{t-1}=\alpha+(\phi-1)y_{t-1}+e_t.
\]

Tedy:
\[
    \Delta y_t=\alpha+\beta_1y_{t-1}+e_t,
\]
kde:
\[
    \beta_1=\phi-1.
\]

Dickey--Fullerův test testuje:
\[
    H_0:\beta_1=0
\]
řada má jednotkový kořen, tedy je nestacionární,

\[
    H_1:\beta_1<0
\]
řada je stacionární.

ADF test, tedy augmented Dickey--Fuller test, přidává zpožděné diference, aby odstranil autokorelaci:
\[
    \Delta y_t =
    \alpha+\beta_1 y_{t-1}+\beta_2t
    +\delta_1\Delta y_{t-1}
    +\cdots+
    \delta_p\Delta y_{t-p}
    +e_t.
\]

Varianty testu:
\begin{itemize}
    \item bez konstanty,
    \item s konstantou,
    \item s konstantou a trendem.
\end{itemize}

Volba varianty závisí na povaze dat. Pokud řada viditelně roste v čase, je vhodné uvažovat variantu
s trendem.

\subsubsection*{1.7 Regrese s časovými řadami a autokorelace chyb}

Základní regresní model pro časové řady:
\[
    y_t=\beta_0+\beta_1x_{t1}+\cdots+\beta_kx_{tk}+u_t.
\]

Maticově:
\[
    y=X\beta+u.
\]

OLS odhad:
\[
    \hat\beta=(X^\top X)^{-1}X^\top y.
\]

Pro časové řady jsou důležité zejména tyto předpoklady:
\begin{itemize}
    \item linearita v parametrech,
    \item žádná perfektní multikolinearita,
    \item striktní exogenita:
    \[
        E(u_t\mid X)=0,
    \]
    \item homoskedasticita:
    \[
        Var(u_t\mid X)=\sigma^2,
    \]
    \item žádná autokorelace chyb:
    \[
        Cor(u_t,u_s\mid X)=0,\qquad t\neq s.
    \]
\end{itemize}

Pokud jsou chyby autokorelované, například:
\[
    u_t=\rho u_{t-1}+e_t,
\]
pak při splnění exogenity může být OLS stále nestranný a konzistentní, ale:
\begin{itemize}
    \item není efektivní,
    \item není BLUE,
    \item běžné směrodatné chyby jsou špatně,
    \item t-testy a F-testy mohou vést k chybným závěrům.
\end{itemize}

\paragraph{Durbin--Watsonův test}

Testuje hlavně AR(1) autokorelaci reziduí:
\[
    DW=
    \frac{\sum_{t=2}^{T}(\hat u_t-\hat u_{t-1})^2}
    {\sum_{t=1}^{T}\hat u_t^2}
    \approx 2(1-\hat\rho).
\]

Interpretace:
\begin{itemize}
    \item $DW\approx 2$: žádná autokorelace,
    \item $DW<2$: pozitivní autokorelace,
    \item $DW>2$: negativní autokorelace.
\end{itemize}

\paragraph{Breusch--Godfreyův test}

Obecnější test autokorelace řádu $p$:
\[
    u_t=\rho_1u_{t-1}+\cdots+\rho_pu_{t-p}+e_t.
\]

Pomocná regrese typicky vysvětluje rezidua $\hat u_t$ pomocí původních regresorů a zpožděných reziduí.

\paragraph{Řešení autokorelace}

\begin{itemize}
    \item \textbf{HAC / Newey--West směrodatné chyby:} neopravují samotný model, ale opravují inferenci.
    \item \textbf{GLS/FGLS:} transformují model tak, aby chyby byly blíže nekorelovaným a homoskedastickým.
    \item \textbf{Cochrane--Orcutt:} používá odhad $\rho$ a transformuje model:
    \[
        \widetilde y_t=y_t-\hat\rho y_{t-1},
    \]
    \[
        \widetilde x_{tj}=x_{tj}-\hat\rho x_{t-1,j}.
    \]
\end{itemize}

\subsubsection*{1.8 Regresní modely se zpožděnými proměnnými}

V ekonomii často účinky nevznikají okamžitě. Například změna úrokové sazby může ovlivnit investice
až s několikaměsíčním zpožděním. Proto používáme modely se zpožděnými proměnnými.

\paragraph{Regrese s ARMA chybami}

Model:
\[
    y_t=\alpha+\beta x_t+u_t,
\]
kde chyba má vlastní dynamiku:
\[
    u_t=
    \sum_{i=1}^{p}\phi_i u_{t-i}
    +
    \sum_{j=1}^{q}\theta_j e_{t-j}
    +e_t.
\]

Tento model říká, že $x_t$ působí současně, ale nevysvětlená část modelu má časovou strukturu.

\paragraph{ARMAX model}

ARMA model s exogenní proměnnou:
\[
    y_t=
    \alpha+\beta x_t
    +
    \sum_{i=1}^{p}\phi_i y_{t-i}
    +
    \sum_{j=1}^{q}\theta_j e_{t-j}
    +e_t.
\]

Rozdíl proti regresi s ARMA chybami je v interpretaci. V ARMAX modelu může $x_t$ působit i nepřímo
přes zpožděné hodnoty $y_t$.

\paragraph{FDL model}

Finite distributed lag model:
\[
    y_t=\alpha+\sum_{j=0}^{q}\delta_j x_{t-j}+u_t.
\]

Koeficient:
\[
    \delta_0
\]
je okamžitý efekt změny $x_t$ na $y_t$.

Dlouhodobý efekt permanentní změny $x$:
\[
    \delta=\sum_{j=0}^{q}\delta_j.
\]

Příklad interpretace: pokud $x$ trvale vzroste o jednu jednotku, potom se po odeznění všech zpoždění
$y$ změní o součet všech $\delta_j$.

\paragraph{ADL model}

Autoregressive distributed lag model:
\[
    y_t=
    \alpha+
    \sum_{i=1}^{p}\rho_i y_{t-i}
    +
    \sum_{j=0}^{q}\delta_j x_{t-j}
    +u_t.
\]

Tento model kombinuje:
\begin{itemize}
    \item zpožděné hodnoty závislé proměnné,
    \item současné i zpožděné hodnoty vysvětlující proměnné.
\end{itemize}

Dlouhodobý multiplikátor v ADL modelu je:
\[
    \frac{\sum_{j=0}^{q}\delta_j}{1-\sum_{i=1}^{p}\rho_i},
\]
pokud je model stabilní.

\subsubsection*{1.9 Predikce a volba modelu}

U časových řad je predikce velmi častým cílem. Predikce může být:
\begin{itemize}
    \item \textbf{bodová}, tedy jedna odhadnutá hodnota,
    \item \textbf{intervalová}, tedy interval možných hodnot.
\end{itemize}

Pro $h$-krokovou predikci označujeme:
\[
    \widehat y_{T+h}.
\]

Při delším predikčním horizontu obvykle roste nejistota, proto se intervaly predikce rozšiřují.

Modely se často porovnávají pomocí:
\[
    AIC=-2\ell+2K,
\]
\[
    BIC=-2\ell+K\ln(T),
\]
kde $\ell$ je log-likelihood a $K$ počet parametrů.

Nižší AIC/BIC znamená lepší kompromis mezi přesností a složitostí modelu.

Pro vyhodnocení predikcí se používají například:
\[
    MAE=\frac{1}{n}\sum_{t=1}^{n}|y_t-\hat y_t|,
\]
\[
    RMSE=\sqrt{\frac{1}{n}\sum_{t=1}^{n}(y_t-\hat y_t)^2}.
\]

Důležité je rozlišovat:
\begin{itemize}
    \item \textbf{in-sample fit}: jak model sedí na datech, na kterých byl odhadnut,
    \item \textbf{out-of-sample predikci}: jak model predikuje nová data.
\end{itemize}

\subsection*{2. Vícerozměrné časové řady, VAR a kointegrace}

\subsubsection*{2.1 Vícerozměrné časové řady}

Pokud sledujeme více časových řad současně, používáme vektor:
\[
    \mathbf y_t=(y_{t1},\dots,y_{tM})^\top.
\]

Například:
\[
    \mathbf y_t =
    \begin{pmatrix}
    HDP_t\\
    inflace_t\\
    urokova\ sazba_t
    \end{pmatrix}.
\]

Vícerozměrné modely jsou důležité, protože ekonomické veličiny se často ovlivňují navzájem.
Například úrokové sazby ovlivňují inflaci, ale zároveň centrální banka reaguje na inflaci úrokovými sazbami.

\subsubsection*{2.2 VAR model}

VAR model, tedy vector autoregression, modeluje více endogenních proměnných najednou.
Každá proměnná je vysvětlována vlastními zpožděními i zpožděními ostatních proměnných.

VAR($p$):
\[
    \mathbf y_t
    =
    \boldsymbol\alpha
    +
    \Phi_1\mathbf y_{t-1}
    +\cdots+
    \Phi_p\mathbf y_{t-p}
    +
    \mathbf e_t.
\]

Pro dvě proměnné a jedno zpoždění:
\[
    y_{1t}=\alpha_1+\phi_{11}y_{1,t-1}+\phi_{12}y_{2,t-1}+e_{1t},
\]
\[
    y_{2t}=\alpha_2+\phi_{21}y_{1,t-1}+\phi_{22}y_{2,t-1}+e_{2t}.
\]

Výhoda VAR modelu:
\begin{itemize}
    \item nemusíme předem určit jednu závislou a jednu nezávislou proměnnou,
    \item všechny proměnné mohou být endogenní,
    \item model zachycuje dynamické vazby mezi proměnnými.
\end{itemize}

Nevýhoda:
\begin{itemize}
    \item rychle roste počet parametrů,
    \item interpretace může být složitá,
    \item model vyžaduje dostatek dat.
\end{itemize}

\paragraph{Redukovaná a strukturální forma}

Redukovaná forma VAR neobsahuje současné hodnoty ostatních endogenních proměnných na pravé straně.
Současné vztahy jsou schovány v korelaci chyb.

Strukturální VAR může mít podobu:
\[
    A_0\mathbf y_t
    =
    \mathbf a
    +
    A_1\mathbf y_{t-1}
    +\cdots+
    A_p\mathbf y_{t-p}
    +
    \boldsymbol\varepsilon_t.
\]

Strukturální forma ale vyžaduje identifikační omezení, protože současné vztahy mezi proměnnými
nelze odhadnout automaticky bez dodatečných předpokladů.

\paragraph{Impulse response function}

Impulse response function popisuje, jak celý systém reaguje na jednorázový šok v jedné proměnné.
Například můžeme zkoumat, jak šok v úrokové sazbě ovlivní inflaci v dalších obdobích.

\subsubsection*{2.3 Grangerova kauzalita}

Proměnná $x$ Grangerovsky způsobuje proměnnou $y$, pokud minulost $x$ pomáhá predikovat $y$
nad rámec minulosti samotné proměnné $y$.

Ve VAR modelu testujeme, zda jsou koeficienty u zpožděných hodnot $x$ v rovnici pro $y$ nulové:
\[
    H_0:\phi_1=\phi_2=\cdots=\phi_p=0.
\]

Pokud $H_0$ zamítneme, minulost $x$ zlepšuje predikci $y$.

Důležité: Grangerova kauzalita není kauzalita v experimentálním smyslu. Znamená predikční vztah,
nikoliv nutně skutečný příčinný efekt.

\subsubsection*{2.4 Integrace a kointegrace}

Řada je integrována řádu $d$, značíme:
\[
    y_t \sim I(d),
\]
pokud je nutné ji $d$-krát diferencovat, aby byla stacionární.

Příklady:
\[
    I(0) = \text{stacionární řada},
\]
\[
    I(1) = \text{řada stacionární po první diferenci}.
\]

Random walk:
\[
    y_t=y_{t-1}+e_t
\]
je typický příklad $I(1)$ procesu.

\paragraph{Kointegrace}

Předpokládejme, že $y_t$ a $x_t$ jsou obě $I(1)$.
Regrese:
\[
    y_t=\alpha+\beta x_t+u_t
\]
může být zdánlivá. Pokud ale reziduum:
\[
    u_t=y_t-\alpha-\beta x_t
\]
je stacionární, tedy $I(0)$, potom jsou $y_t$ a $x_t$ kointegrované.

Intuice: jednotlivé řady mohou dlouhodobě růst nebo bloudit, ale existuje mezi nimi dlouhodobý
rovnovážný vztah, od kterého se neodchylují trvale.

\subsubsection*{2.5 Error Correction Model}

Pokud jsou dvě řady kointegrované, lze použít ECM:
\[
    \Delta y_t =
    \mu+\gamma\Delta x_t
    +\delta(y_{t-1}-\alpha-\beta x_{t-1})
    +e_t.
\]

Interpretace:
\begin{itemize}
    \item $\gamma$ zachycuje krátkodobý efekt změny $x$,
    \item $\alpha+\beta x_{t-1}$ zachycuje dlouhodobou rovnováhu,
    \item výraz $(y_{t-1}-\alpha-\beta x_{t-1})$ je odchylka od rovnováhy,
    \item $\delta$ je rychlost návratu k rovnováze.
\end{itemize}

Obvykle očekáváme:
\[
    \delta<0.
\]

Pokud je $y_{t-1}$ nad dlouhodobou rovnováhou, záporné $\delta$ tlačí $\Delta y_t$ dolů.

\subsubsection*{2.6 VECM}

VECM je vícerozměrná verze ECM pro kointegrované proměnné:
\[
    \Delta\mathbf y_t
    =
    \Pi\mathbf y_{t-1}
    +
    \Gamma_1\Delta\mathbf y_{t-1}
    +\cdots+
    \Gamma_{p-1}\Delta\mathbf y_{t-p+1}
    +
    \mathbf e_t.
\]

Pokud:
\[
    0<rank(\Pi)<M,
\]
existuje kointegrace a můžeme psát:
\[
    \Pi=\gamma\beta^\top.
\]

Zde:
\begin{itemize}
    \item $\beta^\top \mathbf y_t$ jsou kointegrační vztahy,
    \item $\gamma$ zachycuje rychlosti přizpůsobení k rovnováze.
\end{itemize}

\paragraph{Volba mezi VAR a VECM}

\begin{center}
\begin{tabular}{|p{0.42\textwidth}|p{0.45\textwidth}|}
\hline
\textbf{Situace} & \textbf{Vhodný model} \\
\hline
Všechny proměnné jsou $I(0)$ & VAR v úrovních \\
\hline
Proměnné jsou $I(1)$ a nejsou kointegrované & VAR v diferencích \\
\hline
Proměnné jsou $I(1)$ a jsou kointegrované & VECM \\
\hline
Jedna závislá proměnná a vysvětlující proměnné se zpožděními & FDL, ADL, ARMAX nebo ECM \\
\hline
\end{tabular}
\end{center}

\subsection*{3. Panelová data}

\subsubsection*{3.1 Co jsou panelová data}

Panelová data kombinují průřezovou a časovou dimenzi. Sledujeme stejné jednotky v čase:
\[
    y_{it},\qquad i=1,\dots,N,\quad t=1,\dots,T.
\]

Příklady:
\begin{itemize}
    \item stejní lidé sledovaní několik let,
    \item firmy sledované po kvartálech,
    \item země EU sledované po letech,
    \item kraje sledované v čase.
\end{itemize}

Rozdíl mezi pooled daty a panelem:
\begin{itemize}
    \item \textbf{Pooled cross sections:} v různých letech sledujeme různé jednotky.
    \item \textbf{Panelová data:} v různých letech sledujeme stejné jednotky.
\end{itemize}

\subsubsection*{3.2 Typy panelů}

\begin{itemize}
    \item \textbf{Vyvážený panel:} každá jednotka má pozorování ve všech obdobích.
    \item \textbf{Nevyvážený panel:} některá pozorování chybí.
    \item \textbf{Krátký panel:} velké $N$, malé $T$.
    \item \textbf{Dlouhý panel:} malé $N$, velké $T$.
    \item \textbf{Velký panel:} velké $N$ i velké $T$.
\end{itemize}

U krátkých panelů se často používají FE, RE, FD a související metody. U dlouhých panelů se více
objevují problémy podobné časovým řadám, například stacionarita a autokorelace.

\subsubsection*{3.3 Základní panelový model}

Základní panelová regrese:
\[
    y_{it}
    =
    \alpha
    +
    \mathbf x_{it}^\top\beta
    +
    \mathbf z_i^\top\gamma
    +
    u_{it}.
\]

Kde:
\begin{itemize}
    \item $\mathbf x_{it}$ jsou proměnné měnící se v čase,
    \item $\mathbf z_i$ jsou časově neměnné charakteristiky jednotky,
    \item $u_{it}$ je chyba.
\end{itemize}

Často ale některé důležité časově neměnné vlastnosti nepozorujeme. Pak píšeme:
\[
    y_{it}
    =
    \alpha
    +
    \mathbf x_{it}^\top\beta
    +
    c_i
    +
    u_{it}.
\]

Zde:
\[
    c_i
\]
je nepozorovaná heterogenita jednotky.

Příklady $c_i$:
\begin{itemize}
    \item vrozené schopnosti člověka,
    \item firemní kultura,
    \item geografická poloha regionu,
    \item instituce nebo historické faktory země.
\end{itemize}

Hlavní výhoda panelových dat je, že dokážou kontrolovat časově neměnné nepozorované faktory.

\subsubsection*{3.4 Pooled OLS}

Pooled OLS ignoruje panelovou strukturu a odhaduje obyčejnou regresi:
\[
    y_{it}=\alpha+\mathbf x_{it}^\top\beta+v_{it},
\]
kde:
\[
    v_{it}=c_i+u_{it}.
\]

Pooled OLS je problematický, pokud:
\[
    Cor(c_i,\mathbf x_{it})\neq 0.
\]

Pak je $c_i$ vynechaná proměnná korelovaná s regresory a odhad $\hat\beta$ je vychýlený.

Pooled OLS může být použitelný, pokud:
\[
    Cor(c_i,\mathbf X_i)=0,
\]
ale v praxi je to často silný předpoklad.

\subsubsection*{3.5 LSDV}

LSDV, tedy least squares dummy variables, přidává dummy proměnné pro jednotky:
\[
    y_{it}
    =
    \alpha_1I\{i=1\}
    +
    \alpha_2I\{i=2\}
    +\cdots+
    \alpha_NI\{i=N\}
    +
    \mathbf x_{it}^\top\beta
    +
    u_{it}.
\]

Každá jednotka má vlastní intercept:
\[
    \alpha_i=\alpha+c_i.
\]

Výhoda:
\begin{itemize}
    \item explicitně kontroluje nepozorovanou heterogenitu jednotek.
\end{itemize}

Nevýhoda:
\begin{itemize}
    \item při velkém $N$ je mnoho dummy proměnných,
    \item nelze pohodlně odhadovat efekty časově neměnných proměnných.
\end{itemize}

LSDV a FE dávají stejné odhady koeficientů $\beta$.

\subsubsection*{3.6 First Difference estimator}

Model:
\[
    y_{it}
    =
    \alpha
    +
    \mathbf x_{it}^\top\beta
    +
    c_i
    +
    u_{it}.
\]

Pro období $t-1$:
\[
    y_{i,t-1}
    =
    \alpha
    +
    \mathbf x_{i,t-1}^\top\beta
    +
    c_i
    +
    u_{i,t-1}.
\]

Odečtením dostaneme:
\[
    y_{it}-y_{i,t-1}
    =
    (\mathbf x_{it}-\mathbf x_{i,t-1})^\top\beta
    +
    (u_{it}-u_{i,t-1}).
\]

Tedy:
\[
    \Delta y_{it}
    =
    \Delta \mathbf x_{it}^\top\beta
    +
    \Delta u_{it}.
\]

Výhoda:
\[
    c_i
\]
zmizí, protože je v čase konstantní.

Nevýhody:
\begin{itemize}
    \item ztratíme první období,
    \item časově neměnné proměnné nelze odhadnout,
    \item pokud je chyba silně autokorelovaná, může být FD méně vhodný.
\end{itemize}

FD je přirozený zejména pro panel s $T=2$.

\subsubsection*{3.7 Fixed Effects estimator}

FE používá tzv. within transformaci. Nejprve vezmeme průměr přes čas pro jednotku $i$:
\[
    \bar y_i
    =
    \alpha
    +
    \bar{\mathbf x}_i^\top\beta
    +
    c_i
    +
    \bar u_i.
\]

Poté odečteme průměr od původní rovnice:
\[
    y_{it}-\bar y_i
    =
    (\mathbf x_{it}-\bar{\mathbf x}_i)^\top\beta
    +
    (u_{it}-\bar u_i).
\]

Člen $c_i$ zmizí.

FE tedy využívá pouze změny uvnitř jednotky v čase. Neptá se, proč mají některé jednotky vyšší
hodnoty než jiné, ale ptá se, co se stane, když se u téže jednotky změní $x$.

Interpretace:
\[
    \beta_j
\]
říká, jak se změní $y$ při změně $x_j$ uvnitř stejné jednotky v čase, po kontrole všech časově
neměnných vlastností jednotky.

Příklad: Pokud zkoumáme vliv vzdělání na mzdu a používáme individuální FE, kontrolujeme všechny
časově neměnné vlastnosti člověka, například schopnosti nebo rodinné zázemí, pokud se v čase nemění.

Nevýhody FE:
\begin{itemize}
    \item nelze odhadnout efekt časově neměnných proměnných,
    \item využívá jen vnitřní variabilitu,
    \item pokud se proměnná v čase téměř nemění, její efekt se odhaduje nepřesně.
\end{itemize}

\subsubsection*{3.8 Between estimator}

Between estimator používá pouze rozdíly mezi jednotkami. Odhaduje regresi časových průměrů:
\[
    \bar y_i
    =
    \alpha
    +
    \bar{\mathbf x}_i^\top\beta
    +
    e_i.
\]

Nevyužívá změny v čase, pouze to, že některé jednotky mají v průměru vyšší $x$ a jiné nižší $x$.
Proto neumí dobře kontrolovat nepozorovanou heterogenitu korelovanou s regresory.

Obecně je méně důležitý než FE a RE, ale pomáhá pochopit rozdíl mezi:
\begin{itemize}
    \item within variabilitou,
    \item between variabilitou.
\end{itemize}

\subsubsection*{3.9 Random Effects estimator}

Random effects model opět vychází z:
\[
    y_{it}
    =
    \alpha
    +
    \mathbf x_{it}^\top\beta
    +
    c_i
    +
    u_{it}.
\]

Rozdíl proti FE je v předpokladu:
\[
    Cor(c_i,\mathbf X_i)=0.
\]

Tedy nepozorovaná heterogenita jednotky nesmí korelovat s vysvětlujícími proměnnými.

Pokud tento předpoklad platí, RE je efektivnější než FE, protože využívá:
\begin{itemize}
    \item variabilitu uvnitř jednotek,
    \item variabilitu mezi jednotkami.
\end{itemize}

Navíc RE umožňuje odhadovat efekty časově neměnných proměnných.

Složená chyba:
\[
    v_{it}=c_i+u_{it}.
\]

RE používá quasi-demeaning:
\[
    y_{it}-\lambda\bar y_i
    =
    \alpha(1-\lambda)
    +
    (\mathbf x_{it}-\lambda\bar{\mathbf x}_i)^\top\beta
    +
    v_{it}-\lambda\bar v_i.
\]

Kde:
\[
    \lambda
    =
    1
    -
    \sqrt{
    \frac{1}{1+T\frac{Var(c_i)}{Var(u_{it})}}
    }.
\]

Pokud:
\[
    \lambda=0,
\]
RE se podobá pooled OLS.

Pokud:
\[
    \lambda=1,
\]
RE se blíží FE transformaci.

\subsubsection*{3.10 Hausmanův test}

Hausmanův test porovnává FE a RE.

Hypotéza:
\[
    H_0: Cor(c_i,\mathbf X_i)=0.
\]

Pokud platí $H_0$, RE je konzistentní a efektivnější než FE.
Pokud $H_0$ neplatí, RE je nekonzistentní a měli bychom preferovat FE.

Testová statistika:
\[
    H
    =
    (\hat\beta_{FE}-\hat\beta_{RE})^\top
    \left[
    Var(\hat\beta_{FE})-Var(\hat\beta_{RE})
    \right]^{-1}
    (\hat\beta_{FE}-\hat\beta_{RE}).
\]

Rozhodnutí:
\begin{itemize}
    \item nezamítneme $H_0$: RE je přijatelný,
    \item zamítneme $H_0$: preferujeme FE.
\end{itemize}

Praktická intuice:
\begin{itemize}
    \item FE je bezpečnější, pokud se bojíme korelace mezi $c_i$ a regresory,
    \item RE je efektivnější, ale stojí na silnějším předpokladu.
\end{itemize}

\subsubsection*{3.11 Correlated Random Effects}

CRE, tedy correlated random effects nebo Mundlakův přístup, dovoluje, aby nepozorovaná heterogenita
souvisela s průměry regresorů.

Model:
\[
    y_{it}
    =
    \alpha
    +
    \mathbf x_{it}^\top\beta
    +
    \bar{\mathbf x}_i^\top\theta
    +
    a_i
    +
    u_{it}.
\]

Zde:
\[
    \bar{\mathbf x}_i
\]
jsou časové průměry regresorů pro jednotku $i$.

Pokud jsou koeficienty $\theta$ významné, naznačuje to, že běžný RE předpoklad je problematický.
CRE je užitečný kompromis, protože umožňuje:
\begin{itemize}
    \item testovat korelaci mezi heterogenitou a regresory,
    \item interpretovat within i between efekty,
    \item někdy zachovat možnost odhadnout časově neměnné proměnné.
\end{itemize}

\subsubsection*{3.12 Časové efekty a two-way fixed effects}

Někdy existují šoky společné pro všechny jednotky v určitém čase. Například:
\begin{itemize}
    \item ekonomická krize,
    \item pandemie,
    \item změna regulace,
    \item technologická změna.
\end{itemize}

Model s časovými efekty:
\[
    y_{it}
    =
    \alpha
    +
    \mathbf x_{it}^\top\beta
    +
    c_i
    +
    d_t
    +
    u_{it}.
\]

Kde:
\[
    d_t
\]
je efekt společný pro všechny jednotky v čase $t$.

Pokud model obsahuje individuální efekty $c_i$ i časové efekty $d_t$, mluvíme o:
\[
    \text{two-way fixed effects}.
\]

Tento model kontroluje:
\begin{itemize}
    \item všechny časově neměnné rozdíly mezi jednotkami,
    \item všechny šoky společné pro jednotky v daném čase.
\end{itemize}

\subsubsection*{3.13 Robustní inference v panelech}

V panelových modelech často vzniká:
\begin{itemize}
    \item heteroskedasticita mezi jednotkami,
    \item autokorelace uvnitř jednotek,
    \item korelace mezi jednotkami ve stejném čase.
\end{itemize}

Proto se často používají robustní nebo clusterované směrodatné chyby.

\paragraph{Clusterování podle jednotek}

Umožňuje libovolnou korelaci chyb uvnitř jednotky v čase:
\[
    E(u_{it}u_{is})\neq 0.
\]

To je typické, pokud pozorování jedné firmy, osoby nebo země v čase spolu souvisí.

\paragraph{Clusterování podle času}

Umožňuje korelaci mezi jednotkami ve stejném období:
\[
    E(u_{it}u_{jt})\neq 0.
\]

To je důležité při společných makroekonomických šocích.

\paragraph{Two-way clustering}

Kombinuje clusterování podle jednotek i podle času.

Po odhadu panelového modelu je vhodné testovat autokorelaci reziduí, například:
\begin{itemize}
    \item Wooldridgeovým testem pro FE modely,
    \item Breusch--Godfreyho testem pro RE modely.
\end{itemize}

Pokud existuje korelace chyb, lze použít:
\begin{itemize}
    \item robustní/clusterované směrodatné chyby,
    \item FGLS,
    \item vhodnější specifikaci modelu.
\end{itemize}

\subsubsection*{3.14 Dynamické panelové modely}

Dynamický panel obsahuje zpožděnou závisle proměnnou:
\[
    y_{it}
    =
    \alpha
    +
    \rho y_{i,t-1}
    +
    \mathbf x_{it}^\top\beta
    +
    c_i
    +
    u_{it}.
\]

Tento model je vhodný, pokud má závislá proměnná setrvačnost. Například:
\begin{itemize}
    \item současná nezaměstnanost závisí na minulé nezaměstnanosti,
    \item současný zisk firmy závisí na minulém zisku,
    \item současná kriminalita závisí na minulé kriminalitě.
\end{itemize}

Problém: běžný FE odhad je v krátkých panelech vychýlený, protože zpožděná závislá proměnná
koreluje s transformovanou chybou. Tomu se říká dynamická panelová bias, často také Nickell bias.

\paragraph{Arellano--Bond estimator}

Začneme modelem:
\[
    y_{it}
    =
    \alpha
    +
    \rho y_{i,t-1}
    +
    \mathbf x_{it}^\top\beta
    +
    c_i
    +
    u_{it}.
\]

Nejprve diferencujeme, čímž odstraníme $c_i$:
\[
    \Delta y_{it}
    =
    \rho \Delta y_{i,t-1}
    +
    \Delta \mathbf x_{it}^\top\beta
    +
    \Delta u_{it}.
\]

Ale:
\[
    \Delta y_{i,t-1}=y_{i,t-1}-y_{i,t-2}
\]
je korelováno s:
\[
    \Delta u_{it}=u_{it}-u_{i,t-1}.
\]

Proto používáme instrumenty, typicky hlubší zpoždění:
\[
    y_{i,t-2},\ y_{i,t-3},\dots
\]

Odhad se provádí pomocí GMM.

Důležité diagnostiky:
\begin{itemize}
    \item test autokorelace v diferencovaných reziduích,
    \item u Arellano--Bond odhadu očekáváme AR(1) v diferencích,
    \item neměli bychom mít AR(2) v diferencích,
    \item používá se také Sarganův nebo Hansenův test instrumentů.
\end{itemize}

\subsection*{4. Shrnutí rozdílů mezi hlavními modely}

\begin{center}
\begin{tabular}{|p{0.36\textwidth}|p{0.54\textwidth}|}
\hline
\textbf{Situace} & \textbf{Vhodný přístup} \\
\hline
Jedna stacionární časová řada & AR, MA, ARMA \\
\hline
Nestacionární řada, která se po diferencování stane stacionární & ARIMA \\
\hline
Sezónní časová řada & Sezónní dummy nebo SARIMA \\
\hline
Regrese časových řad s autokorelací chyb & HAC SE, GLS/FGLS, ARMA chyby \\
\hline
Efekt proměnné se zpožděním & FDL, ADL, ARMAX \\
\hline
Více vzájemně propojených endogenních časových řad & VAR \\
\hline
Nestacionární, ale kointegrované řady & ECM nebo VECM \\
\hline
Panel s nepozorovanou heterogenitou korelovanou s regresory & FE nebo FD \\
\hline
Panel s heterogenitou nekorelovanou s regresory & RE \\
\hline
Volba mezi FE a RE & Hausmanův test a věcná argumentace \\
\hline
Panel s šoky společnými pro všechna období & časové efekty nebo two-way FE \\
\hline
Dynamický panel se zpožděným $y$ & Arellano--Bond / GMM \\
\hline
\end{tabular}
\end{center}

\subsection*{5. Nejdůležitější rovnice k zapamatování}

\[
    \text{AR(p):}\qquad
    y_t=\alpha+\sum_{i=1}^{p}\phi_i y_{t-i}+e_t.
\]

\[
    \text{ARMA(p,q):}\qquad
    y_t=
    \alpha+
    \sum_{i=1}^{p}\phi_i y_{t-i}
    +
    \sum_{j=1}^{q}\theta_j e_{t-j}
    +e_t.
\]

\[
    \text{ARIMA(p,d,q):}\qquad
    \Delta^d y_t=
    \alpha+
    \sum_{i=1}^{p}\phi_i\Delta^d y_{t-i}
    +
    \sum_{j=1}^{q}\theta_j e_{t-j}
    +e_t.
\]

\[
    \text{ADF test:}\qquad
    \Delta y_t =
    \alpha+\beta_1 y_{t-1}+\beta_2t
    +\delta_1\Delta y_{t-1}
    +\cdots+
    \delta_p\Delta y_{t-p}
    +e_t.
\]

\[
    \text{FDL:}\qquad
    y_t=\alpha+\sum_{j=0}^{q}\delta_jx_{t-j}+u_t.
\]

\[
    \text{ADL:}\qquad
    y_t=
    \alpha+
    \sum_{i=1}^{p}\rho_i y_{t-i}
    +
    \sum_{j=0}^{q}\delta_jx_{t-j}
    +u_t.
\]

\[
    \text{VAR(p):}\qquad
    \mathbf y_t
    =
    \boldsymbol\alpha
    +
    \sum_{i=1}^{p}\Phi_i\mathbf y_{t-i}
    +
    \mathbf e_t.
\]

\[
    \text{ECM:}\qquad
    \Delta y_t =
    \mu+\gamma\Delta x_t
    +\delta(y_{t-1}-\alpha-\beta x_{t-1})
    +e_t.
\]

\[
    \text{Panel FE:}\qquad
    y_{it}-\bar y_i
    =
    (\mathbf x_{it}-\bar{\mathbf x}_i)^\top\beta
    +
    (u_{it}-\bar u_i).
\]

\[
    \text{Panel FD:}\qquad
    \Delta y_{it}
    =
    \Delta \mathbf x_{it}^\top\beta
    +
    \Delta u_{it}.
\]

\[
    \text{Panel RE:}\qquad
    y_{it}-\lambda\bar y_i
    =
    \alpha(1-\lambda)
    +
    (\mathbf x_{it}-\lambda\bar{\mathbf x}_i)^\top\beta
    +
    v_{it}-\lambda\bar v_i.
\]

\[
    \text{Dynamický panel:}\qquad
    y_{it}
    =
    \alpha
    +
    \rho y_{i,t-1}
    +
    \mathbf x_{it}^\top\beta
    +
    c_i
    +
    u_{it}.
\]


\newpage
